\documentclass[UTF8]{ctexart}
\usepackage[a4paper,margin=2.4cm]{geometry}
\usepackage{amsmath,amssymb,amsthm,mathtools}
\usepackage{enumitem}
\usepackage{hyperref}

\title{论文证明树笔记：<paper title>}
\author{Codex 阅读笔记}
\date{\today}

\newcommand{\claim}[1]{\paragraph{命题：}#1}
\newcommand{\role}[1]{\paragraph{作用：}#1}
\newcommand{\deps}[1]{\paragraph{依赖：}#1}
\newcommand{\source}[1]{\paragraph{来源：}#1}

\begin{document}
\maketitle

\section{论文信息}
\begin{itemize}[leftmargin=2em]
  \item 标题：<title>
  \item 作者：<authors>
  \item arXiv：<arxiv id>
  \item PDF：\texttt{<pdf path>}
  \item 源码目录：\texttt{<source dir>}
\end{itemize}

\section{主定理}
\claim{<用自己的话陈述主结论。>}
\role{<说明该结论解决的问题，以及它相比已有工作的推进。>}

\section{证明树总览}
\begin{enumerate}[label=\textbf{T\arabic*.}, leftmargin=2.5em]
  \item <主定理>
  \begin{enumerate}[label*=\textbf{.\arabic*}, leftmargin=2.5em]
    \item <关键归约>
    \item <核心技术定理>
    \item <从核心定理推出主定理>
  \end{enumerate}
\end{enumerate}

\section{节点详解}

\subsection{节点 1：<名称>}
\source{\texttt{<tex file>}，标签 \texttt{<label>}，或章节 <section>.}
\claim{<该节点的数学断言。>}
\deps{<它依赖哪些定义、引理、定理或外部文献。>}
\role{<它在整棵证明树中的位置。>}

\section{关键公式与估计}
\begin{align*}
  <formula>
\end{align*}

\section{主要创新}
\begin{itemize}[leftmargin=2em]
  \item <创新点 1>
  \item <创新点 2>
\end{itemize}

\section{待核对}
\begin{itemize}[leftmargin=2em]
  \item <不确定点、需要补读的引理或外部文献。>
\end{itemize}

\end{document}
